\section*{الفصل \ref{ch:b2:ism-star-formation} --- الوسط بين النجمي وتكوّن النجوم}

\begin{solution}{exo:b2:ism-star-formation:1}
$n = \num{2.7e-24}/\num{1.67e-24} = 1.6$ ذرّة هيدروجين في السنتيمتر المكعب. وفي الهواء
$n = P/k_{\mathrm B}T = 101325/(\num{1.381e-23}\times288) = \qty{2.55e25}{m^{-3}}$،
أي \num{2.5e19} جزيء في السنتيمتر المكعب: أكثر بنحو $1.6\times10^{19}$
مرة.
\end{solution}

\begin{solution}{exo:b2:ism-star-formation:2}
$M_J \propto n^{-1/2}$: قسمة $n$ على 100 تضرب $M_J$ في 10، فتكون
$M_J = 54.6$ كتلة شمسية عند $10^2$ جسيم في السنتيمتر المكعب. وتحوي السحابة
$\num{24000}/54.6 \approx 440$ \omterm{def:b2:ism-star-formation:jeans}{كتلة جينز}: فلا يمكن أن تكون مستقرة ككلّ، ولا بدّ أن
تتشظّى.
\end{solution}

\begin{solution}{exo:b2:ism-star-formation:3}
$c_s = \sqrt{\num{1.381e-23}\times10/(2.33\times\num{1.661e-27})} =
\qty{189}{m.s^{-1}}$ عند \qty{10}{K}؛ وأكبر بعامل $\sqrt{10}$، أي \qty{597}{m.s^{-1}}،
عند \qty{100}{K}. ولعبور \qty{0.1}{pc}: $\num{3.086e15}/189 = \qty{1.63e13}{s}$،
أي \qty{0.52}{Myr}.
\end{solution}

\begin{solution}{exo:b2:ism-star-formation:4}
$\rho = 2.33\times\num{1.661e-27}\times n$ حيث $n$ بوحدة \unit{m^{-3}}. فمن أجل $10^5$
جسيم في السنتيمتر المكعب، $\rho = \qty{3.87e-16}{kg.m^{-3}}$
و$t_{\mathrm{ff}} = \sqrt{3\pi/(32G\rho)} = \qty{3.38e12}{s} = \qty{0.107}{Myr}$؛ ومن أجل
$10^2$ تكون $\rho$ أصغر بعامل \num{1000} و$t_{\mathrm{ff}}$ أطول بعامل $\sqrt{1000} = 31.6$،
 أي \qty{3.38}{Myr}. \omterm{def:b2:ism-star-formation:molecular-cloud}{فاللبّ الكثيف} ينهار أسرع بنحو 32 مرة.
\end{solution}

\begin{solution}{exo:b2:ism-star-formation:5}
بما أن $k_{\mathrm B}T/(\mu m_u) = c_s^2$، فإن $(5k_{\mathrm B}T/G\mu m_u)^{3/2} =
5^{3/2}c_s^3G^{-3/2}$، ومنه الصيغة. وبما أن $M_J \propto \rho^{-1/2}$،
فإن خفض $M_J$ إلى النصف يتطلب كثافة أعلى بأربع مرات. وبما أن $M_J \propto T^{3/2}$،
فإن التسخين من \qty{10}{K} إلى \qty{20}{K} يضربها في $2^{3/2} = 2.83$.
\end{solution}

\begin{solution}{exo:b2:ism-star-formation:6}
$t_{\mathrm{KH}} = GM^2/(RL) = \num{6.674e-11}\times(\num{1.989e30})^2/(2\times
\num{6.96e8}\times2\times\num{3.83e26}) = \qty{2.48e14}{s}$، أي
\qty{7.8}{Myr}. وهو أطول من عمر معظم الأقراص (يزول نصفها بعد
نحو \qty{3}{Myr}): فالنجم لا يزال ينكمش حين يختفي قرصه.
\end{solution}

\begin{solution}{exo:b2:ism-star-formation:7}
عند \qty{10}{K}، $c_s^3/G = 189^3/\num{6.674e-11} = \qty{1.01e17}{kg.s^{-1}}$، أي
$\num{1.01e17}\times\num{3.156e7}/\num{1.989e30} = \num{1.6e-6}$ كتلة شمسية في
السنة، وتستغرق الكتلة الشمسية الواحدة \qty{0.62}{Myr}. وعند \qty{20}{K} يكون المعدل
أعلى بعامل $2^{3/2} = 2.83$، أي \num{4.5e-6} كتلة شمسية في السنة، وتستغرق الكتلة
الشمسية الواحدة \qty{0.22}{Myr}.
\end{solution}

\begin{solution}{exo:b2:ism-star-formation:8}
$r = \sqrt{L/(16\pi\sigma)}/T^2$ مع $\sqrt{\num{3.83e26}/(16\pi\times\num{5.67e-8})}
= \qty{1.159e16}{m.K^2}$: عند \qty{1300}{K}، $r = \qty{6.86e9}{m} = \qty{0.046}{au}$،
أي داخل مدار عطارد بكثير. وعند \qty{5.2}{au}، $T = 278/\sqrt{5.2} = \qty{122}{K}$.
\end{solution}

\begin{solution}{exo:b2:ism-star-formation:9}
$v = \sqrt{GM/r} = \sqrt{\num{6.674e-11}\times1.3\times\num{1.989e30}/\num{1.496e13}}
= \qty{3.40}{km.s^{-1}}$، و$P = 2\pi r/v = \qty{2.77e10}{s} = 877$ سنة
(أو، على نحو مكافئ، $\sqrt{100^3/1.3}$ سنة). وفي مليون سنة يُتمّ الغاز
$10^6/877 \approx \num{1140}$ دورة: فالقرص قديم مقارنة بدورانه الذاتي،
ولهذا تبدو حلقاته دوائر منتظمة.
\end{solution}

\begin{solution}{exo:b2:ism-star-formation:10}
$M = \int_0^R 4\pi r^2\rho_0R/r\,\mathrm dr = 2\pi\rho_0R^3$، وداخل $r$
$m(r) = 2\pi\rho_0Rr^2$. ومع $\mathrm dm = 4\pi\rho_0Rr\,\mathrm dr$،
\[
  U = -\int_0^R \frac{G\,m(r)}{r}\,\mathrm dm
    = -8\pi^2G\rho_0^2R^2\int_0^R r^2\,\mathrm dr = -\frac{8\pi^2}{3}G\rho_0^2R^5
    = -\frac23\,\frac{GM^2}{R},
\]
باستعمال $\rho_0 = M/(2\pi R^3)$. وهي $\tfrac{2/3}{3/5} = 1.11$ ضعف قيمة الكرة المتجانسة:
فالسحابة المركّزة أشدّ ارتباطًا وتحتاج إلى طاقة حرارية أكثر بنسبة \qty{11}{\percent}
لإسنادها؛ فهي تنهار بسهولة أكبر.
\end{solution}

\begin{solution}{exo:b2:ism-star-formation:11}
الدور $2\pi/\Omega = \qty{6.28e14}{s}$، أي نحو \qty{20}{Myr}. ومع
$R = \qty{3.086e15}{m}$ يكون $j = \Omega R^2 = \qty{9.52e16}{m^2.s^{-1}}$،
و$r_c = j^2/(GM) = \num{9.07e33}/\num{1.327e20} = \qty{6.83e13}{m} = \qty{457}{au}$.
وفي الدوران الصلب ينمو $j = \Omega r^2$ مع مربع البُعد عن المحور:
فالغاز الآتي من داخل اللبّ أقل زخمًا زاويًا، فيستقر أقرب إلى النجم.
\end{solution}

\begin{solution}{exo:b2:ism-star-formation:12}
$\bar t = 2.46$، $\bar f = 0.588$؛ $\sum(t - \bar t)^2 = 12.612$،
$\sum(t - \bar t)(f - \bar f) = -2.018$، ومنه $b = 2.018/12.612 = 0.160$ لكل مليون
سنة و$a = 0.588 + 0.160\times2.46 = 0.982$. يبلغ المستقيم $0.5$ عند
$(0.982 - 0.5)/0.160 = \qty{3.0}{Myr}$ والصفر عند $0.982/0.160 = \qty{6.1}{Myr}$، كما في
الشكل. ولا يمكن لنسبة أن تنزل تحت الصفر: فالعناقيد الأقدم، مثل NGC 1960 عند
\qty{30}{Myr}، تحتفظ بنسبة صغيرة قرب الصفر لا يستطيع المستقيم وصفها.
\end{solution}

\begin{solution}{pb:b2:ism-star-formation:1}
\textbf{1.} $R = 0.08\times\num{3.086e16} = \qty{2.469e15}{m}$.

\textbf{2.} $N = \qty{1.5e26}{m^{-2}}$، ومنه $2.8\times\num{1.661e-27}\times\num{1.5e26} =
\qty{0.698}{kg.m^{-2}}$.

\textbf{3.} $M = 0.698\times\pi R^2 = 0.698\times\num{1.915e31} = \qty{1.34e31}{kg}$،
أي $6.7$ كتلة شمسية.

\textbf{4.} $V = \tfrac43\pi R^3 = \qty{6.30e46}{m^3}$ و
$\rho = M/V = \qty{2.12e-16}{kg.m^{-3}}$.

\textbf{5.} $n = \rho/(2.33m_u) = \qty{5.48e10}{m^{-3}}$، أي \num{5.5e4} جسيم
في السنتيمتر المكعب.

\textbf{6.} نعم: كثافته بين $10^4$ و$10^5$، ونصف قطره بين 0.03
و\qty{0.2}{pc}؛ أمّا كتلته، 6.7 كتلة شمسية، فعند الحدّ الأعلى لمدى اللبوب.

\textbf{7.} $M_J = 5.46\times(10^4/\num{5.48e4})^{1/2} = 2.33$ كتلة شمسية.

\textbf{8.} $6.7/2.33 = 2.9$ \omterm{def:b2:ism-star-formation:jeans}{كتلة جينز}: فاللبّ غير مستقر وينهار.

\textbf{9.} $R_J = 0.1316\times(10^4/\num{5.48e4})^{1/2} = \qty{0.056}{pc}$، وهو أصغر من
نصف القطر، \qty{0.08}{pc}: والخلاصة نفسها.

\textbf{10.} $2^{3/2}\times2.33 = 6.6$ كتلة شمسية: أي كتلة اللبّ تقريبًا. فعند
\qty{20}{K} يكون اللبّ على الحدّ؛ فالبرودة هي ما يجعله ينهار.

\textbf{11.} $t_{\mathrm{ff}} = \sqrt{3\pi/(32\times\num{6.674e-11}\times\num{2.12e-16})}
= \qty{4.56e12}{s}$، أي نحو \num{145000} سنة.

\textbf{12.} $R/c_s = \num{2.469e15}/189 = \qty{1.31e13}{s}$، أي \num{414000} سنة: أطول بثلاث
مرات من السقوط الحر. فلا تستطيع موجة ضغط عبور اللبّ قبل أن
ينهار؛ ولا يستطيع الضغط أن يستجيب بسرعة كافية لإيقاف الانهيار.

\textbf{13.} $c_s = \qty{189}{m.s^{-1}}$ (\cref{exo:b2:ism-star-formation:3}).

\textbf{14.} $c_s^3/G = \qty{1.01e17}{kg.s^{-1}}$، أي \num{1.6e-6} كتلة شمسية في السنة.

\textbf{15.} $1/\num{1.6e-6} = \num{624000}$ سنة.

\textbf{16.} ثلث 6.7 هو $2.2$: نجمان كتلة كلٍّ منهما كتلة شمسية واحدة.

\textbf{17.} $t_{\mathrm{KH}} = GM_\odot^2/(3R_\odot\times4L_\odot) = \qty{8.25e13}{s}$،
أي \qty{2.6}{Myr}.

\textbf{18.} أطول، أربعة أضعاف زمن المراكمة: فحين تتوقف المراكمة يكون النجم
لا يزال ينكمش ويسطع دون اندماج؛ إنه \omterm{def:b2:ism-star-formation:t-tauri}{نجم تي الثور}، لم يبلغ النسق
الرئيسي بعد.

\textbf{19.} $\Omega = \qty{0.4e3}{m.s^{-1}}/\qty{3.086e16}{m} = \qty{1.30e-14}{rad.s^{-1}}$.

\textbf{20.} $j = 1.30\times10^{-14}\times(\num{2.469e15})^2 = \qty{7.90e16}{m^2.s^{-1}}$.

\textbf{21.} $r_c = j^2/(GM_\odot) = \num{6.24e33}/\num{1.327e20} = \qty{4.70e13}{m}$، أي
\qty{314}{au}: أكبر من قرص إتش إل الثور، كما هو متوقع لغاز آتٍ من حافة
اللبّ.

\textbf{22.} $T = (L/16\pi\sigma r^2)^{1/4} = \qty{278}{K}$ عند \qty{1}{au}،
و$278/\sqrt{30} = \qty{51}{K}$ عند \qty{30}{au}.

\textbf{23.} $r = \num{1.159e16}/180^2 = \qty{3.58e11}{m}$، أي \qty{2.4}{au}، كما في النص.

\textbf{24.} يستغرق النجم ذو الكتلة الشمسية الواحدة
\textbf{\qty[text-series-to-math=true]{624}{} ألف سنة} ليتجمّع، أي نحو أربعة
أضعاف \omterm{def:b2:ism-star-formation:free-fall}{زمن السقوط الحر} للبّه (\num{145000} سنة): فالغاز الخارجي لا يصل إلى
\omterm{def:b2:ism-star-formation:protostar}{النجم الأوّلي} إلا تدريجيًا، من الداخل إلى الخارج.
\end{solution}
